\documentclass[11pt]{article}

\usepackage{amsmath,amssymb,amsthm,mathtools}
\usepackage[margin=1in]{geometry}
\usepackage{microtype}

\newtheorem{theorem}{Theorem}[section]
\newtheorem{lemma}[theorem]{Lemma}
\newtheorem{corollary}[theorem]{Corollary}
\newtheorem{remark}[theorem]{Remark}

\newcommand{\dd}{\,\mathrm d}
\newcommand{\ip}[2]{\left\langle #1,#2\right\rangle}

\title{Triangles with Mixed Rooted Branches:\\
A Degree-Biased Jensen Bound}
\author{}
\date{}

\begin{document}

\maketitle

\begin{abstract}
For integers $r,s\geq0$, let $R_{r,s}$ be a triangle with a distinguished
vertex $o$, with $r$ pendant edges and $s$ internally disjoint two-edge
paths attached at $o$.  We prove that every graphon of edge density
$p\geq1/2$ satisfies
\[
 t(R_{r,s},W)\geq p^{r+2s+1}(2p-1)=\Phi_{R_{r,s}}(p).
\]
The cases $s=0$ and $r=0$ were known separately.  The new point is that the
two branch types may occur simultaneously.  A complement-kernel
symmetrization proves that the mean two-edge-tail density under the measure
biased by $d^r$ is at least $p^2$.  A single convex Jensen step then treats
all $s$ tails at once.  In the six-vertex catalogue this proves the formerly
open Atlas graph $95$.
\end{abstract}


%%%%%%%%%%%%%%%%%%%%%%%%%%%%%%%%%%%%%%%%%%%%%%%%%%%%%%%%%%%%%%%%%%%%%%%%%%%%%%%
\section{Definitions and the target}
%%%%%%%%%%%%%%%%%%%%%%%%%%%%%%%%%%%%%%%%%%%%%%%%%%%%%%%%%%%%%%%%%%%%%%%%%%%%%%%

Let $(\Omega,\mu)$ be an arbitrary probability space, and let
$W\colon\Omega^2\to[0,1]$ be a symmetric measurable graphon.  Write
\[
 p=\int_{\Omega^2}W(x,y)\dd\mu(x)\dd\mu(y),
 \qquad
 d(x)=\int_\Omega W(x,y)\dd\mu(y),
\]
and put
\[
 A(x)=(T_Wd)(x)=\int_\Omega W(x,y)d(y)\dd\mu(y).
\]
Thus $A(x)$ is the rooted density of a two-edge path whose initial vertex
is fixed at $x$.  Define the rooted triangle density
\[
 \tau(x)=\int_{\Omega^2}
 W(x,y)W(x,z)W(y,z)\dd\mu(y)\dd\mu(z).
\]

For $r,s\geq0$, construct $R_{r,s}$ from a triangle rooted at $o$ by
adjoining $r$ distinct leaves adjacent to $o$, and $s$ pairwise
vertex-disjoint paths $o-u_i-v_i$.  No other vertices are identified and no
other edges are added.  Integrating the branch vertices first gives the
exact identity
\begin{equation}
 t(R_{r,s},W)=\int_\Omega \tau(x)d(x)^rA(x)^s\dd\mu(x).
 \label{eq:density-identity}
\end{equation}

The graph has $3+r+2s$ vertices and the same number of edges.  Once the
triangle is properly coloured, every branch vertex has $x-1$ choices after
its parent is coloured.  Hence
\[
 \chi_{R_{r,s}}(x)=x(x-1)^{r+2s+1}(x-2),
 \qquad \chi(R_{r,s})=3.
\]
For $q=1-p>0$, direct substitution gives
\begin{align}
 \Phi_{R_{r,s}}(p)
 &=q^{3+r+2s}\frac1q
   \left(\frac pq\right)^{r+2s+1}
   \left(\frac{2p-1}{q}\right)
 \notag\\
 &=p^{r+2s+1}(2p-1).
 \label{eq:target}
\end{align}
This polynomial also gives the continued value $1$ at $p=1$.


%%%%%%%%%%%%%%%%%%%%%%%%%%%%%%%%%%%%%%%%%%%%%%%%%%%%%%%%%%%%%%%%%%%%%%%%%%%%%%%
\section{The weighted tail mean}
%%%%%%%%%%%%%%%%%%%%%%%%%%%%%%%%%%%%%%%%%%%%%%%%%%%%%%%%%%%%%%%%%%%%%%%%%%%%%%%

Put $M_j=\int_\Omega d(x)^j\dd\mu(x)$ for $j\geq0$.  The next lemma is the
ingredient that permits the two branch types to be mixed.

\begin{lemma}[Degree-biased two-edge-tail mean]
\label{lem:biased-tail}
For every integer $r\geq0$ and every graphon of edge density
$p\in[1/2,1]$,
\[
 \int_\Omega d(x)^rA(x)\dd\mu(x)\geq p^2M_r.
\]
Equivalently, under the probability measure
\[
 \dd\nu_r(x)=\frac{d(x)^r}{M_r}\dd\mu(x),
\]
one has $\mathbb E_{\nu_r}A\geq p^2$.
\end{lemma}

\begin{proof}
The denominator $M_r$ is positive because $M_r\geq p^r>0$ for $r\geq1$
by Jensen, while $M_0=1$.  Let $U=1-W$, let $T_U$ be its integral operator,
and put $u=T_U\mathbf1=1-d$ and $q=1-p$.  Since
\[
 T_Ud=T_U(\mathbf1-u)=u-T_Uu,
\]
we have
\[
 A=T_Wd=p\mathbf1-T_Ud=d-q\mathbf1+T_Uu.
\]
Consequently, if $J_r=\int d^rA$, then
\begin{equation}
 J_r=M_{r+1}-qM_r+\ip{d^r}{T_Uu}.
 \label{eq:J-decomposition}
\end{equation}

Symmetry of $U$ gives
\begin{align*}
 2\left(\ip{d^r}{T_Uu}-\int d^ru^2\right)
 ={}&\int_{\Omega^2}U(x,y)
 \bigl(u(y)-u(x)\bigr)
 \bigl(d(x)^r-d(y)^r\bigr)\dd\mu^2.
\end{align*}
The integrand is nonnegative: $d=1-u$, so $u$ and $d^r$ are oppositely
ordered.  Thus \eqref{eq:J-decomposition} implies
\begin{equation}
 J_r\geq M_{r+2}-M_{r+1}+pM_r.
 \label{eq:J-moments}
\end{equation}

Set $m=M_{r+1}/M_r$.  Cauchy--Schwarz gives
$M_{r+2}/M_r\geq m^2$.  Also $m\geq p$: for $r=0$ this is equality, while
for $r\geq1$ it follows from
\[
 M_{r+1}-pM_r
 =\frac12\int_{\Omega^2}
 (d(x)^r-d(y)^r)(d(x)-d(y))\dd\mu(x)\dd\mu(y)\geq0.
\]
Since $p\geq1/2$ in the application, $m\geq p\geq1-p$, and therefore
\[
 \frac{J_r}{M_r}-p^2
 \geq m^2-m+p-p^2
 =(m-p)(m-(1-p))\geq0.
\]
This proves the assertion.
\end{proof}


%%%%%%%%%%%%%%%%%%%%%%%%%%%%%%%%%%%%%%%%%%%%%%%%%%%%%%%%%%%%%%%%%%%%%%%%%%%%%%%
\section{Rooted triangle inputs}
%%%%%%%%%%%%%%%%%%%%%%%%%%%%%%%%%%%%%%%%%%%%%%%%%%%%%%%%%%%%%%%%%%%%%%%%%%%%%%%

We record the two elementary estimates needed for the main proof.

\begin{lemma}[Pointwise rooted triangle]
\label{lem:pointwise-triangle}
For almost every $x$,
\[
 \tau(x)\geq\max\{2A(x)-p,0\}.
\]
\end{lemma}

\begin{proof}
The elementary inequality $ab\geq a+b-1$ for $a,b\in[0,1]$, applied to
$W(x,y)$ and $W(x,z)$ and multiplied by $W(y,z)\geq0$, gives after
integration
\[
 \tau(x)\geq2\int_\Omega W(x,y)d(y)\dd\mu(y)-p=2A(x)-p.
\]
Also $\tau(x)\geq0$.
\end{proof}

\begin{lemma}[Weighted rooted triangle]
\label{lem:weighted-triangle}
For every integer $r\geq0$ and every graphon of edge density $p>0$,
\[
 \int_\Omega d(x)^r\tau(x)\dd\mu(x)
 \geq\frac{2p-1}{p}M_{r+2}.
\]
\end{lemma}

\begin{proof}
Put $I_r=\int d^r\tau$ and $J_r=\int d^rA$.  Applying
$ab\geq a+b-1$ to the two edges incident with the fixed root, multiplying
by $d(x)^rW(y,z)$, and integrating yields
\[
 I_r\geq2J_r-pM_r.
\]
Using \eqref{eq:J-moments},
\[
 I_r\geq2M_{r+2}-2M_{r+1}+pM_r.
\]
Subtracting $(2p-1)M_{r+2}/p$ from the right-hand side leaves
\[
 \frac1p\left(M_{r+2}-2pM_{r+1}+p^2M_r\right)
 =\frac1p\int_\Omega d(x)^r(d(x)-p)^2\dd\mu(x)\geq0.
\]
\end{proof}

\begin{lemma}[Tail convexity]
\label{lem:tail-convexity}
For $p>0$ and every integer $s\geq1$, the function
\[
 f_{p,s}(a)=a^s\max\{2a-p,0\},\qquad 0\leq a\leq1,
\]
is convex and nondecreasing.
\end{lemma}

\begin{proof}
It vanishes on $[0,p/2]$.  On $[p/2,1]$ it is
$2a^{s+1}-pa^s$, whose first derivative is
$a^{s-1}(2(s+1)a-ps)>0$ and whose second derivative is
$sa^{s-2}(2(s+1)a-p(s-1))\geq0$.  At $p/2$ the right derivative is
nonnegative and no smaller than the zero left derivative.  This proves both
claims on the whole interval.
\end{proof}


%%%%%%%%%%%%%%%%%%%%%%%%%%%%%%%%%%%%%%%%%%%%%%%%%%%%%%%%%%%%%%%%%%%%%%%%%%%%%%%
\section{The mixed-branch theorem}
%%%%%%%%%%%%%%%%%%%%%%%%%%%%%%%%%%%%%%%%%%%%%%%%%%%%%%%%%%%%%%%%%%%%%%%%%%%%%%%

\begin{theorem}[Mixed rooted triangle branches]
\label{thm:main}
For all integers $r,s\geq0$ and every graphon $W$ whose edge density lies in
$p\in[1/2,1]$,
\[
 \boxed{
 t(R_{r,s},W)
 \geq p^{r+2s+1}(2p-1)
 =\Phi_{R_{r,s}}(p).
 }
\]
This is the complete interval required by $\chi(R_{r,s})=3$.
\end{theorem}

\begin{proof}
First suppose $s=0$.  By \eqref{eq:density-identity},
Lemma~\ref{lem:weighted-triangle}, and Jensen,
\[
 t(R_{r,0},W)
 \geq\frac{2p-1}{p}M_{r+2}
 \geq p^{r+1}(2p-1).
\]

Now suppose $s\geq1$.  Under the probability measure $\nu_r$ from
Lemma~\ref{lem:biased-tail}, \eqref{eq:density-identity} and
Lemma~\ref{lem:pointwise-triangle} give
\[
 t(R_{r,s},W)\geq M_r\,\mathbb E_{\nu_r} f_{p,s}(A).
\]
By Lemmas~\ref{lem:tail-convexity} and \ref{lem:biased-tail},
\[
 \mathbb E_{\nu_r} f_{p,s}(A)
 \geq f_{p,s}(\mathbb E_{\nu_r}A)
 \geq f_{p,s}(p^2).
\]
Since $p\geq1/2$, one has $p^2\geq p/2$, and hence
\[
 f_{p,s}(p^2)=p^{2s}(2p^2-p)=p^{2s+1}(2p-1).
\]
Finally $M_r\geq p^r$ by Jensen (including the identity $M_0=1$), so
\[
 t(R_{r,s},W)\geq p^{r+2s+1}(2p-1),
\]
which equals \eqref{eq:target}.
\end{proof}

At $p=1/2$ the target vanishes and every step remains valid.  At $p=1$,
$0\leq W\leq1$ and $\int W=1$ imply $W=1$ almost everywhere, so the density
and target both equal $1$.  No division by $2p-1$ occurs.


%%%%%%%%%%%%%%%%%%%%%%%%%%%%%%%%%%%%%%%%%%%%%%%%%%%%%%%%%%%%%%%%%%%%%%%%%%%%%%%
\section{Catalogue consequence}
%%%%%%%%%%%%%%%%%%%%%%%%%%%%%%%%%%%%%%%%%%%%%%%%%%%%%%%%%%%%%%%%%%%%%%%%%%%%%%%

\begin{corollary}[Atlas 95]
NetworkX Atlas graph $95$, with graph6 code \texttt{EG\char125?}, is
positive.  It has order and size six, chromatic number three, and
\[
 \chi_{F_{95}}(x)=x(x-1)^4(x-2).
\]
For every graphon $W$ of edge density $p\in[1/2,1]$,
\[
 t(F_{95},W)\geq p^4(2p-1)=\Phi_{F_{95}}(p).
\]
\end{corollary}

\begin{proof}
Atlas $95$ has NetworkX edge list
\[
 04,\ 05,\ 12,\ 14,\ 24,\ 34.
\]
The vertices $1,2,4$ form a triangle.  Root it at $4$: vertex $3$ is a
direct leaf, while $4-0-5$ is a two-edge path, and there are no other edges.
Thus the graph is $R_{1,1}$, and Theorem~\ref{thm:main} applies.  The
chromatic polynomial and target follow from the formulas in Section~1.
\end{proof}


%%%%%%%%%%%%%%%%%%%%%%%%%%%%%%%%%%%%%%%%%%%%%%%%%%%%%%%%%%%%%%%%%%%%%%%%%%%%%%%
\section{Formalization boundary and exact limitations}
%%%%%%%%%%%%%%%%%%%%%%%%%%%%%%%%%%%%%%%%%%%%%%%%%%%%%%%%%%%%%%%%%%%%%%%%%%%%%%%

The proof uses ordinary, non-injective graphon homomorphism densities on an
arbitrary probability space.  All integrands are bounded and measurable, so
Tonelli--Fubini applies.  The only analytic inputs are Jensen,
Cauchy--Schwarz, and the displayed opposite-order symmetrization for $U$.
No finite-step reduction, regularity assumption, minimum-degree hypothesis,
or injective-copy interpretation is present.

For a Lean formalization, the new reusable statement is
Lemma~\ref{lem:biased-tail}.  Its clean signature should include
$p\geq1/2$ explicitly.  The remaining ingredients already occur in the
rooted triangle--tree development: the density factorization, pointwise
rooted-triangle bound, tail convexity, weighted rooted-triangle inequality,
and moment Jensen bound.  The graph-specific layer need only exhibit Atlas
$95\cong R_{1,1}$ and transport the generic result.

The theorem covers only a bouquet at one triangle root whose branches have
length one or two and are internally disjoint.  It does not cover a
three-edge path, two length-two paths sharing their intermediate vertex,
branches at different triangle vertices, or attachments to a larger clique
or triangle book.  In particular it does not classify Atlas $97$, $100$, or
$102$.  The false stronger assertion
$t(F+\text{leaf},W)\geq p\,t(F,W)$ is explicitly refuted in
\texttt{FAILED\_PROOF\_ATTEMPTS.md}; it is not used here.


%%%%%%%%%%%%%%%%%%%%%%%%%%%%%%%%%%%%%%%%%%%%%%%%%%%%%%%%%%%%%%%%%%%%%%%%%%%%%%%
\section{What the Lean formalization actually proves}
%%%%%%%%%%%%%%%%%%%%%%%%%%%%%%%%%%%%%%%%%%%%%%%%%%%%%%%%%%%%%%%%%%%%%%%%%%%%%%%

The machine-checked development is
\texttt{lean/Taeyoung/Methods/MixedBranch.lean}, ending at
\texttt{satisfiesLowerBound\_95}.  It proves the single case
$(r,s)=(1,1)$ --- the only one of six vertices --- and it does so by an
argument that is \emph{not} the one above.

\paragraph{The departure: no symmetrization, no Jensen.}
This note handles general $(r,s)$ by a complement-kernel symmetrization,
showing that the mean two-edge-tail density under the $d^{r}$-biased measure is
at least $p^{2}$, and then disposes of all $s$ tails with one convex Jensen
step.  At $(1,1)$ neither is needed.  Peeling gives
\[
 t(R_{1,1},W)=\int\tau\,A\,d,\qquad A=T_Wd,
\]
and the pointwise Goodman bound $\tau\geq2A-p$, weighted by $Ad\geq0$, already
reduces the problem to two moments of the single weight $d\cdot A$:
\[
 t\geq2C-pB,\qquad B=\int dA,\quad C=\int dA^{2}.
\]
The weighted Cauchy--Schwarz $B^{2}\leq pC$ and path Sidorenko $B\geq p^{3}$
then close it through
\begin{equation}
 p\,(2C-pB)-p\,(2p^{5}-p^{4})\;\geq\;(B-p^{3})\bigl(2(B+p^{3})-p^{2}\bigr)\;\geq\;0,
 \label{eq:lean-mixed}
\end{equation}
both factors being nonnegative because $B\geq p^{3}>0$ and $2p^{3}\geq p^{2}$ on
$p\geq1/2$.  Every input predates the row:
\texttt{Link.rootedTriangle\_ge},
\texttt{PureChordal.integral\_mul\_sq\_le\_integral\_mul\_integral\_mul\_sq}
with weight $d$, and \texttt{PathSidorenko.pow\_three\_le\_pathIntegral}.

So the two proofs agree on $(1,1)$ and diverge in generality: this note's is
the one that scales in $r$ and $s$, while \eqref{eq:lean-mixed} is specific to
one tail and one leaf.  If a seven-vertex member ever enters scope, the
symmetrization is the thing to formalize.

\paragraph{Everything else transcribes directly.}
The peeling is `Methods/Peeling.lean`'s generic ladder plus three collapses;
$\tau$ appears as its own definition, so the triangle needs no separate lemma.

\end{document}
